% Motivationsschreiben DI Alexander Sommer - Bewerbung OeNB
% Compiler: pdfLaTeX / Tectonic
\documentclass[11pt]{article}
\usepackage[a4paper,top=20mm,bottom=18mm,left=24mm,right=24mm]{geometry}
\usepackage[T1]{fontenc}
\usepackage[utf8]{inputenc}
\usepackage[ngerman]{babel}
\IfFileExists{sourcesanspro.sty}%
  {\usepackage[default,semibold]{sourcesanspro}}%
  {\usepackage[scaled=0.92]{helvet}\renewcommand{\familydefault}{\sfdefault}}
\usepackage{microtype}
\usepackage{xcolor}
\usepackage[hidelinks]{hyperref}

\definecolor{accent}{HTML}{1F3A5F}
\definecolor{muted}{HTML}{5A6472}

\hypersetup{pdftitle={Motivationsschreiben DI Alexander Sommer},pdfauthor={DI Alexander Sommer}}
\setlength{\parindent}{0pt}
\setlength{\parskip}{7pt}
\linespread{1.06}
\pagestyle{empty}

\begin{document}

% ---------------------------------------------------------------- Absender / Empfänger
\begin{minipage}[t]{0.55\textwidth}
  \textbf{\color{accent}DI Alexander Sommer}\\
  Gregorygasse 12/6/2\\
  1230 Wien\\[4pt]
  +43 681 102 104 17\\
  \href{mailto:alexander.sommer@gmail.com}{alexander.sommer@gmail.com}
\end{minipage}\hfill
\begin{minipage}[t]{0.40\textwidth}\raggedleft
  Oesterreichische Nationalbank\\
  Personalmanagement\\
  Otto-Wagner-Platz 3\\
  1090 Wien
\end{minipage}

\vspace{10pt}
\hfill Wien, \today

\vspace{6pt}
{\bfseries\color{accent}Bewerbung als Leiter der Abteilung IT Workplace Services}

\vspace{2pt}
Sehr geehrte Damen und Herren,

vier Millionen Euro jährliches Einsparpotenzial strukturiert identifiziert. Bis zu 40.000 Euro pro Jahr durch konsequente Lizenz-Governance eingespart. Eine geschäftskritische Applikationslandschaft mit mehr als 85 Systemen in einem der sensibelsten Infrastruktursektoren Österreichs verantwortet. Das ist die Bilanz, die ich in mehr als 15 Jahren IT-Führung bei OMV, Wiener Wohnen und zuletzt der Austrian Power Grid AG erarbeitet habe -- aktuell als Teamleiter Digital Workplace Solutions, einer Funktion, die sich mit der ausgeschriebenen Leitung von IT Workplace Services fast deckungsgleich beschreiben lässt. Aktuell treibe ich dort unter anderem die Neuaufstellung des Dokumentenmanagements (DMS\,2.0) und die Neugestaltung des Intranets voran, beides auf Basis von Microsoft\,365, und hinterfrage dabei bewusst auch unsere M365-Governance im Kontext digitaler Souveränität. Mit dieser Erfahrung bewerbe ich mich nun bei der Oesterreichischen Nationalbank: einer Institution, die wie kaum eine andere für Stabilität, Verlässlichkeit und volkswirtschaftliche Verantwortung steht und deren Aufgabe mich seit Langem fasziniert.

Mein Interesse an Finanzmärkten ist dabei keine Momentaufnahme: Vor einigen Jahren habe ich gemeinsam mit drei Partnern eine automatisierte Handelsplattform für den Live-Handel an NYSE und NASDAQ entwickelt -- über die auch das Wall Street Journal berichtete (Ausgabe vom 10. August 2015). Auch wenn dieses Kapitel mittlerweile zurückliegt, ist mir das Zusammenspiel aus Technologie, Risikosteuerung und Marktdynamik seither vertraut geblieben und prägt bis heute meinen Blick auf IT in regulierten, risikosensiblen Umfeldern.

Führung verstehe ich als Gestaltungsaufgabe, nicht als Verwaltung: Im Rahmen des konzernweiten Programms Change4Growth habe ich mein früheres Team -- verantwortlich für eine Applikationslandschaft von mehr als 340 Systemen -- bewusst neu aufgestellt und näher an die strategische Entscheidungsebene gerückt; die entstandene Schwesterstruktur steuere ich bis heute in einer Matrixfunktion mit, ebenso wie eine Vielzahl externer Dienstleister -- eine Konstellation, die der in Ihrer Ausschreibung beschriebenen Steuerung des ausgelagerten First-Level-Supports sehr nahekommt. Diese Führungsqualität wurde 2023 auch extern bestätigt: In einer Potenzialeinschätzung (PEP, IMDE) zählte mein Führungsprofil zu den oberen 10\,\% der Benchmarkgruppe vergleichbarer Führungskräfte. Gerne nenne ich Ihnen auf Wunsch Referenzen, die dies aus eigener Erfahrung bestätigen können -- unter anderem Markus Kasinger (Leiter des Bundesamtes für Cybersicherheit) sowie Bernhard Mittermann, der unter meiner Führung begonnen hat und heute als CIO bei Rexel tätig ist.

Neue Technologien beurteile ich dabei lieber selbst, als mich auf Trends zu verlassen: Neben eigenen, praktischen Projekten mit Large Language Models und Agentic AI nutze ich seit Jahren auch meinen Zugang als Gartner-Seat-Holder (zuletzt Technical-, davor Executive Membership), um Entwicklungen unabhängig einzuordnen, bevor ich sie in die Modernisierungsvorhaben meines Teams einbringe.

Über die Gelegenheit zu einem persönlichen Gespräch würde ich mich sehr freuen.

\vspace{6pt}
Mit freundlichen Grüßen\\[14pt]
DI Alexander Sommer

\end{document}
